\documentclass[a4paper,11pt]{article}

\usepackage{revy}
\usepackage[utf8]{inputenc}
\usepackage[T1]{fontenc}
\usepackage[danish]{babel}
\usepackage{amsmath,amssymb}

\revyname{MatRevy}
\revyyear{2026}
\version{2026-10-09}
\eta{1 minutter}

\title{Pingu}
\author{Philip '23}

\begin{document}
\maketitle

\begin{roles}
\role{P}[Skuespiller] Pingu
\role{F}[Skuespiller] Dum fysiker
\end{roles}

\begin{sketch}
\scene{Vi befinder os nede i sydenden. Pingu skal lave et epsilon-delta bevis. Pingu finder en tavle og struggler med at samle et kridt op i sine luffer. Det lykkedes endelig at få fat i kridtet og Pingu begynder at skrive definitionen på kontinuitet op. Fysikeren begynder på tavlen ved siden af at løse en opgave om heat equation hvor de betragter en pingvins varmetab til omgivelserne over tid. Pingu skæver interesseret over når det går op for ham at fysikerens opgave handler om pingviner. Han vender blikket tilbage til sin egen tavle. Fysikeren lader for enkelthedens skyld i sin opgave en pingvin være cylindrisk. (Måske tegner de en cylinder på tavlen og skriver pingvin over). Pingu arbejder videre i uvidenhed og fløjter muntert for sig selv indtil han kigger over og ser fysikerens arbejde. Han stopper brat op og holder op med at fløjte. Han kigger mistroisk fra "pingvinen" på tavlen ned på sit cylindriske kridt, ned på sig selv lidt frem og tilbage. Han løfter langsomt blikket ud på publikum.}

\says{P}[Vredt og højlydt] NOOT NOOT!

\scene{Lys ned}
\end{sketch}
\end{document}
